% ---------------------------------------------------------------------------
% A unifying LENS (not a theorem): one elementary variance identity organises the
% five empirical settings. Empirical paper -- this is framing, the experiments are
% the contribution. Drop into the criterion section.
% ---------------------------------------------------------------------------
\subsection{A unifying view: the missing signal is within-group potential variance}
\label{sec:unifying}

We do not claim a theorem, and the contribution of this paper is empirical. But a single
back-of-the-envelope identity makes all five settings fall out of \emph{one} quantity, and
we use it only as an organising lens so the reader leaves with one picture rather than five.

Group-relative methods turn a group of $G$ rollouts into advantages by centring each reward
on the group mean, $A_i = R_i - \bar R$. Write each reward as outcome plus a weighted
potential, $R_i = O_i + \beta\,\Phi_i$, with $O_i\in\{0,1\}$ the terminal outcome and
$\Phi_i$ the verifiable process summary (goals discharged, stages completed, fraction of
tests passing). How much the group can move the policy is carried by the \emph{spread} of
its advantages, and the spread of a sum is elementary:
\[
  \operatorname{Var}_G(R)\;=\;\operatorname{Var}_G(O)\;+\;\beta^2\operatorname{Var}_G(\Phi)
  \;+\;2\beta\,\operatorname{Cov}_G(O,\Phi).
\]
That is the entire observation. Its value is not that it is true---it is true by the algebra
of variances---but that each term is one of the criterion's gates, and each term is exactly
what one of our experiments isolates:

\begin{itemize}[leftmargin=1.4em,itemsep=2pt,topsep=2pt]
\item \textbf{Outcome term $\to$ all-fail blindness.} When every rollout fails, $O_i\equiv 0$,
this term and the cross term drop out, and with no potential ($\beta=0$) the group has zero
spread: the dead update. This is the regime the theorem-proving efficiency result lives in.
\item \textbf{Potential term $\to$ finer \& reachable.} The process channel adds spread only
when $\operatorname{Var}_G(\Phi)>0$: $\Phi$ must take more than one value (finer than the
binary outcome) and the policy must actually reach those values (reachable). A potential
equal to the outcome, or one never moved off zero, adds nothing---precisely the software-repair
null, where $\Phi$ stays at $0$ on all $156$ rollouts.
\item \textbf{Cross term $\to$ un-gameable.} Spread alone does not say the update points
\emph{toward} success; the covariance does. A signal whose spread is uncorrelated with, or
pulls away from, the outcome optimises something else---the compliance attractor a misaligned
penalty falls into.
\end{itemize}

We are deliberate about what this is and is not. It is an organising identity that explains
why the same three properties keep deciding the result across very different domains; it is
\emph{not} a proof that process rewards help, and the cross-term reading in particular is a
statement about training \emph{dynamics} that only the experiments establish (\S\ref{sec:law}).
We offer it so that a practitioner can carry one sentence into a new domain:
\textbf{a useful dense process reward is the within-group variance the outcome cannot
supply---present when that variance is non-zero (finer, reachable) and useful when it is
aligned (un-gameable)}---and then check those properties on a handful of base rollouts before
spending any compute.
